% Neutral JAMA-inspired Chinese template, checked against a supplied issue layout.
% No publisher logo, real paper text, authors, DOI, figures, data or download marks.
\documentclass[UTF8,fontset=none]{ctexart}
\usepackage[paperwidth=612bp,paperheight=792bp,left=72.312bp,right=72.72bp,top=67bp,bottom=51bp,headheight=18bp,headsep=14bp,footskip=27bp]{geometry}
% Shared dependencies and neutral material. Compile with XeLaTeX.
% No publisher logos, paper text, data, figures, or download metadata are included.
\usepackage{amsmath,amssymb,graphicx,xcolor,tikz,fancyhdr,needspace,xurl,hyperref,balance,etoolbox}
% The legacy balance package omits the first split box's final baseline depth.
% Reserve that measured depth rather than hiding a CJK overfull-box warning.
\makeatletter
\patchcmd{\@BAlancecol}{\setbox\@leftcolumn=\vbox to \@tempdima}{\advance\@tempdima by \dp1\relax\setbox\@leftcolumn=\vbox to \@tempdima}{}{\PackageWarning{cns-template}{Balance depth patch unavailable; inspect final columns}}
\makeatother
\usetikzlibrary{arrows.meta,positioning}
\usepackage[protrusion=true,expansion=false]{microtype}
\IfFontExistsTF{Noto Serif CJK SC}{%
  \setCJKmainfont{Noto Serif CJK SC}%
}{%
  \setCJKmainfont[BoldFont=FandolSong-Bold.otf,ItalicFont=FandolKai-Regular.otf]{FandolSong-Regular.otf}%
}
\IfFontExistsTF{Noto Sans CJK SC}{%
  \setCJKsansfont{Noto Sans CJK SC}%
}{%
  \setCJKsansfont[BoldFont=FandolHei-Bold.otf]{FandolHei-Regular.otf}%
}
\setmainfont{TeX Gyre Termes}\setsansfont{TeX Gyre Heros}
\setmonofont{TeX Gyre Cursor}
\definecolor{JournalBlue}{HTML}{0097C9}
\definecolor{JournalRed}{HTML}{B83443}
\definecolor{ScienceRed}{HTML}{CD1021}
\definecolor{CiteBlue}{HTML}{0000EE}
\hypersetup{colorlinks=true,linkcolor=JournalBlue,urlcolor=JournalBlue,
  pdftitle={CNS论文汉化：中性版式示例},pdfauthor={示例模板；非期刊官方模板},
  pdfsubject={Neutral source-inspired Chinese layout demonstration; no real research claims}}
\urlstyle{same}\raggedbottom
\setlength{\parindent}{1em}\setlength{\parskip}{0pt}
\setlength{\emergencystretch}{2em}\tolerance=2200
\XeTeXlinebreaklocale "zh"\XeTeXlinebreakskip=0pt plus .1pt
\pagestyle{fancy}\fancyhf{}
\renewcommand{\headrulewidth}{0pt}\renewcommand{\footrulewidth}{0pt}
\newcommand{\PaperTitle}{中文题名占位：从组织结构到证据解释}
\newcommand{\SampleAuthors}{Author A\textsuperscript{1}, Author B\textsuperscript{2,*}}
\newcommand{\SampleAffiliations}{\textsuperscript{1}作者单位一（示例占位）。\quad
  \textsuperscript{2}作者单位二（示例占位）。\par
  *通信作者信息待填。DOI、收稿和出版日期均应从目标原文核实。}
\newcommand{\SampleAbstract}{本段是中性摘要占位文字，不陈述任何真实研究发现。它演示中文题名、作者信息、摘要和双栏正文之间的层级关系。实际使用时，应完整替换为目标论文的忠实中文译文，并保留研究对象、方法、结果和结论的限定条件。图注、参考文献、作者声明和来源信息均属于独立的覆盖项。模板中的图形仅说明排版结构，不表示生物学机制、实验结果或统计关系。}
\newcommand{\DemoCaption}{\textbf{图1. 中性证据链示意（占位图，无实验数据）。}
  \textbf{A}，输入节点表示需要核对的原文材料。
  \textbf{B}，处理节点表示逐段转换及复核。
  \textbf{C}，输出节点表示正文、图注和书目等独立交付部分。
  箭头仅表示本示例的阅读顺序，不代表因果关系、效应方向或数量变化。
  该图由模板自行绘制，所有节点均为文字占位；无样本量、误差线、统计检验或比例尺。
  实际论文图注应逐项保留原稿报告的样本量、单位、检验方法、限定及图形署名。}
\newcommand{\DemoVector}{%
 \resizebox{\linewidth}{!}{%
 \begin{tikzpicture}[font=\sffamily\small,node distance=18mm,
   box/.style={draw=black!70,rounded corners=2pt,fill=black!3,minimum width=26mm,minimum height=16mm,align=center},
   flow/.style={-{Stealth[length=2.5mm]},line width=.7pt}]
  \node[box] (a) {\textbf{A}\quad 输入\\原文材料};
  \node[box,right=of a] (b) {\textbf{B}\quad 处理\\转换与复核};
  \node[box,right=of b] (c) {\textbf{C}\quad 输出\\中文译稿};
  \draw[flow] (a)--(b);\draw[flow] (b)--(c);
  \node[below=11mm of a,align=center] {标题与摘要\\正文与图表};
  \node[below=11mm of b,align=center] {保留限定\\检查单位};
  \node[below=11mm of c,align=center] {可追溯引文\\明确原文信息};
  \node[below=29mm of b,font=\sffamily\footnotesize,text=black!65] {示意占位；不包含真实实验数据};
 \end{tikzpicture}}%
}
\newcommand{\SampleBibliography}{%
 \section*{参考文献（示例占位）}
 {\fontsize{7.6bp}{10bp}\selectfont\setlength{\parindent}{0pt}\interlinepenalty=10000
 \hypertarget{ref:1}{}1. Author A. (20XX). Example reference title. \textit{Journal placeholder}. 卷、页码与DOI待填。此条仅供演示引文跳转。\par\smallskip
 \hypertarget{ref:2}{}2. Author B. (20YY). Another example reference title. \textit{Journal placeholder}. 出版信息待填；不要把占位条目带入正式译稿。\par}
}
\newcommand{\SampleSourceInfo}{%
 \section*{原文信息}
 以下字段全部为示例占位：作者；英文原标题；期刊；年份；卷期；页码或文章号；DOI或来源链接。实际使用时必须逐项替换并核对。\par
 本文件是\JournalName 风格的非官方中文译稿样板，依据一次已核验文章版式进行结构化改编，不代表该期刊所有年代或文体的统一标准，也不保证像素级复刻。无真实论文内容或期刊标识图像。\par
}
% Deliberately neutral demo prose, not translated from a published paper.
\newcommand{\SampleOpening}{%
 \section*{引言：示例正文与引文}
 这是一段用于检查双栏阅读节奏的中性示例文字。段落中的术语、缩写和标点需要在实际项目中与源文保持一致。中文译稿应呈现原论文的完整论述，而不是把多段证据压缩成一段摘要。这里的引文只链接到本文件末尾明确标为占位的参考文献\SampleCiteOne。\par
 若原句仅表示相关、可能或提示，译文也应保留相同的证据强度。数字、单位、正负号、上下标和范围端点应逐项核对。原文存在内部矛盾时，应分别保留相应位置的陈述，再用简短译注说明；不能根据常识把不同的数字静默统一。\par
 \subsection*{内容层级与跨页衔接}
 此小标题演示二级层级。实际论文可能采用不同的小标题密度，模板不会替代对原文结构的识别。双栏文章需按左栏、右栏和跨页续句的顺序读取。若段落在页末中断，应在完成整句核对后接入下一页，避免把图中文字误当正文，或重复翻译同一段。\par
 图表具有独立的信息范围。正文引用到某幅图，并不意味着其图注可以省略。图注中的每个分面、比较对象、时间点、样本量及统计说明都应对应源文。英文图内标签可以保留，并在图旁或文末提供中文释义；实际采用哪种安排，需要依据图形尺寸和译文长度判断\SampleCiteTwo。\par
 \subsection*{来源与可维护数据}
 本段继续演示正文密度。可编辑文本和图注数据应与最终版面分开维护，以便修订时准确定位变化。若只调整版式，应能证明正文、图注和书目内容没有被意外改写。更换字体或纸张后，须重新检查断行、图注与图片的相对位置以及页脚可读性。\par
 文内引用与书目之间应有明确对应关系。这里的两条演示书目没有真实DOI，也没有虚构的实验信息。使用模板时可增删条目，但必须同时更新编号、链接目标和引用位置。页码只表示中文译稿自己的分页，不用于冒充原出版页码。\par
 \subsection*{结果叙述的演示占位}
 本节的文字仅用于呈现长段落在双栏页面中的流动，不报告真实结果。实际翻译应保留观察对象、实验条件和结局指标之间的关系。例如，原文分别描述体外实验与动物实验时，不能把两类结果合并为同一层面的结论；研究使用的对照、时间顺序和分析范围也必须保留。图表中的比较范围不应扩大为对整个人群或全部疾病阶段的判断。\par
 论述中的转折与限制同样属于正文。若研究先提出一种解释，再讨论另一种可能，译文应沿用这一推进顺序。作者明确写出的不确定性应能在中文中被识别；如果尚未确定某个现象的方向或机制，就不能改成已经建立的因果关系。对于术语近似但研究层级不同的概念，可在首次出现时保留缩写，后文保持统一。\par
 \subsection*{方法信息与可复核性}
 方法部分往往包含密集数字和专有名词。翻译时须区分浓度、总量、频率和处理时长，不能因为中文句式调整而改变参数的归属。试剂名称、软件版本和数据访问标识通常保留原样；上下限、排除规则和显著性标准则需要逐字符核对。这里没有提供任何真实实验参数，避免把模板占位误认为可以直接执行的方法。\par
 表格与补充材料的引用也需要保留原层级。正文对补充图的指向并不意味着模板中已经附有该补充图；若源文件只包含主文，应如实标注所交付的范围。图表编号、附录顺序和参考文献数量都应由目标文章决定，本示例中的结构不构成固定配额，也不要求所有论文采用同样的章节。\par
}
\newcommand{\SampleClosing}{%
 \section*{讨论：版式适配与范围}
 上方矢量图是模板原创的排版占位图，不来自任何论文。图形与完整演示图注被放在同一页，便于读者直接核对分面说明。真实项目应使用经过授权或用户提供的原图素材，并保留图注署名与许可信息。尺寸不足时，可适度缩放或调整栏宽，不以删除图注解决空间问题。\par
 模板仅记录一次已核验文章的可观察设计特征。即使文章出自同一期刊，不同年代、研究论文与综述、校正清样与正式版本也可能采用不同的页眉和引文形式。应以用户提供的实际PDF为准，不能仅凭期刊名称推断配色、字号或引用方式\SampleCiteOne。\par
 \subsection*{演示检查项}
 本样板用于展示可编译性、中文字形、标题层级、跨栏图注、页眉页脚和内部链接。图中没有实际数据，因此不应被当作科学插图引用。替换内容后，还需检查所有页的渲染结果，确认无截字、碰撞、异常空白和遗漏段落。\par
 \SampleBibliography
 \SampleSourceInfo
}


\usepackage{multicol,tabularx,colortbl,titlesec}
\newcommand{\JournalName}{JAMA}
\linespread{1}
\definecolor{JAMARed}{HTML}{D71635}\definecolor{JAMABlack}{HTML}{231F20}
\definecolor{JAMABanner}{rgb}{.5947966576,.1058670953,.1181658655}
\definecolor{JAMABeige}{rgb}{.9454795122,.9452658892,.9072098732}
\color{JAMABlack}\setlength{\columnsep}{15.5bp}
\hypersetup{linkcolor=JAMABlack,urlcolor=JAMARed,pdftitle={JAMA风格中文中性样板},pdfauthor={中性示例；非期刊官方模板}}
\titleformat{\section}{\sffamily\bfseries\fontsize{11bp}{13bp}\selectfont}{}{0pt}{}[{\vspace{2bp}\titlerule[.3bp]}]
\titlespacing*{\section}{0pt}{10bp}{6bp}
\titleformat{\subsection}{\sffamily\bfseries\color{JAMARed}\fontsize{9bp}{11bp}\selectfont}{}{0pt}{}
\titlespacing*{\subsection}{0pt}{8bp}{3bp}
\fancyhead[L]{\sffamily\fontsize{7.8bp}{9bp}\selectfont 原创研究\enspace|\enspace 中性中文排版示例}
\fancyhead[R]{\sffamily\fontsize{7.8bp}{9bp}\selectfont Research}
\fancyfoot[L]{\sffamily\fontsize{6.6bp}{8bp}\selectfont 非官方中文样板｜仅供学术交流；出版物版权归原作者及出版商所有。}
\fancyfoot[R]{\sffamily\fontsize{7bp}{8bp}\selectfont\thepage}
\fancypagestyle{first}{\fancyhead{}}
\AddToHook{shipout/background}{\begin{tikzpicture}[remember picture,overlay]\ifnum\value{page}=1\fill[JAMABanner]([xshift=72.312bp,yshift=-11.933bp]current page.north west)rectangle([xshift=539.28bp,yshift=-47.854bp]current page.north west);\draw[white,line width=.4bp]([xshift=126bp,yshift=-11.933bp]current page.north west)--([xshift=126bp,yshift=-47.854bp]current page.north west);\node[anchor=west,text=white,font=\sffamily\bfseries\fontsize{8bp}{10bp}\selectfont]at([xshift=79bp,yshift=-30bp]current page.north west){Research};\else\fill[JAMABanner]([xshift=72.312bp,yshift=-12bp]current page.north west)rectangle([xshift=539.28bp,yshift=-24bp]current page.north west);\fi\end{tikzpicture}}
\newcommand{\absrow}[2]{\noindent\textcolor{JAMARed}{\sffamily\bfseries #1}\enspace #2\par\smallskip\textcolor{white}{\rule{\linewidth}{.35bp}}\par\smallskip}
\newcommand{\neutralpara}{本段为中性排版占位，不陈述真实研究结果。实际译稿须按原文完整替换，保留研究范围、方法条件、统计含义及不确定性。图表与正文应交替安排，使读者在首次讨论结果附近即可查阅对应证据。段落长度、行宽和图文位置可随中文内容调整，不应为填满版面删去限定或虚构数据。}
\begin{document}\thispagestyle{first}\fontsize{8.8bp}{11.6bp}\selectfont
\noindent{\sffamily\bfseries\fontsize{8.5bp}{10bp}\selectfont JAMA\enspace|\enspace\textcolor{JAMARed}{原创研究}}\par\vspace{5bp}
{\sffamily\bfseries\fontsize{14.5bp}{17.2bp}\selectfont\noindent 中文题名占位：研究问题、评价方法与证据解释在目标人群中的应用及其限定条件\par}
\vspace{3bp}{\sffamily\fontsize{13bp}{16bp}\selectfont\noindent 随机临床试验文体样板\par}
\vspace{9bp}{\sffamily\fontsize{6.8bp}{8.8bp}\selectfont\noindent Author A, Degree；Author B, Degree；Author C, Degree\par}
\vspace{12bp}\noindent\begin{minipage}[t]{347.2bp}\vspace{0pt}\setlength{\fboxsep}{11bp}\colorbox{JAMABeige}{\begin{minipage}{325.2bp}\sffamily\fontsize{8.8bp}{11.7bp}\selectfont\setlength{\parindent}{0pt}
\absrow{重要性}{\neutralpara}
\absrow{目的}{说明原文所要回答的研究问题、适用人群与比较对象。此处仅演示结构化摘要的红色标签与黑色正文，不构成临床建议。}
\absrow{设计、场所与参与者}{按目标原文填写研究设计、中心与地区、研究时间、入排标准及参与者构成。对观察性研究、试验或综述，应使用其实际报告的摘要标签。}
\absrow{干预}{填写干预、对照、给药或观察程序及随访长度，保留剂量、单位、分配比例和治疗条件。未报告的信息不得推断补齐。}
\absrow{主要结局与测量}{说明预先规定的主要及次要结局，区分估计目标、观测值、插补值与模型估计，不将统计相关改写成因果效应。}
\absrow{结果}{\neutralpara\par 所有样本量、百分比、置信区间及显著性指标应与原文逐项核对。此中性样板仅保留叙述位置，不提供任何模拟数值，也不放置真实研究图。}
\absrow{结论}{结论应忠实保留适用范围与限制。版式还原不能替代医学审阅；文献翻译不等于个体诊疗建议。}
\absrow{试验注册}{如原文有注册平台与编号，逐项填写并核对链接；此处为占位。}
{\fontsize{6.6bp}{8.3bp}\selectfont 年份、卷期、页码及DOI均待按目标原文填写。\par 在线发表与更正信息不得由文件名推断。\par}
\end{minipage}}\end{minipage}\hfill\begin{minipage}[t]{108bp}\vspace{0pt}\sffamily\fontsize{7.1bp}{9.5bp}\selectfont
\textcolor{JAMARed}{\textbf{+}}\enspace 相关资源位置\par\medskip
\textcolor{JAMARed}{\textbf{←}}\enspace 同期内容位置\par\medskip
\textcolor{JAMARed}{\textbf{+}}\enspace 多媒体位置\par\medskip
\textcolor{JAMARed}{\textbf{←}}\enspace 相关文章位置\par\medskip
\textcolor{JAMARed}{\textbf{+}}\enspace 补充材料位置\par\medskip
\textcolor{JAMARed}{\textbf{+}}\enspace 教育资源位置\par
\vspace{374bp}\textbf{作者单位：}示例单位占位，正式稿应完整保留。\par\medskip
\textbf{通信作者：}姓名、地址、电子邮箱均按原文填写。\par\medskip
\end{minipage}
\clearpage
\noindent\colorbox{JAMABeige}{\begin{minipage}{\dimexpr\linewidth-2\fboxsep}\sffamily\fontsize{8.5bp}{11.1bp}\selectfont\textbf{要点}\par\smallskip\textbf{问题}\quad 在此填写研究问题。\par\smallskip\textbf{发现}\quad 在此填写主要发现，并保留原文的证据层级。\par\smallskip\textbf{意义}\quad 在此填写适用范围及临床或学术含义。\end{minipage}}\par
\begin{multicols}{2}\fontsize{8.8bp}{11.6bp}\selectfont
\section*{背景}\neutralpara\par\neutralpara
\section*{方法}\subsection*{研究设计}\neutralpara\par\subsection*{结局与统计}\neutralpara
\section*{结果}\neutralpara\par 本文采用黑色上标引文\textsuperscript{\hyperlink{ref:1}{1}}；图表保留独立标题、缩写说明及完整脚注。原刊使用细小无衬线标签，中文样板适度扩大正文字号以兼顾阅读。
\end{multicols}
\vspace{5bp}{\sffamily\fontsize{8bp}{10.2bp}\selectfont\noindent\textbf{表1．可编辑中文表格的中性结构}\par\smallskip
\setlength{\tabcolsep}{5bp}\renewcommand{\arraystretch}{1.1}\noindent\begin{tabularx}{\linewidth}{@{}Xll@{}}\hline\rowcolor{JAMABeige}字段&组别A&组别B\\\hline
研究人群&待核对&待核对\\主要结局&待核对&待核对\\次要结局&待核对&待核对\\不良事件&待核对&待核对\\\hline\end{tabularx}\par\smallskip
{\fontsize{7.6bp}{9.2bp}\selectfont 脚注：此表仅演示文字与行带样式，无研究数值；真实中文表须完整保留原表全部列、单位、统计值与脚注。}\par}
\begin{multicols}{2}\fontsize{8.8bp}{11.6bp}\selectfont\subsection*{图表位置}\neutralpara\par\subsection*{可复用范围}此样板依据一次已核验的原创研究版式制作。其他年份、文章类型或显著版式变体须重新查看目标PDF，不能机械套用。
\end{multicols}
\clearpage
\noindent\begin{minipage}[t]{347bp}\vspace{0pt}\DemoVector\end{minipage}\hfill\begin{minipage}[t]{108bp}\vspace{0pt}\sffamily\fontsize{7.6bp}{9.2bp}\selectfont\DemoCaption\end{minipage}\par\vspace{12bp}
\begin{multicols}{2}\fontsize{8.8bp}{11.6bp}\selectfont\section*{讨论}\neutralpara\par\neutralpara
\subsection*{局限性}正式译稿应完整保留局限性，并区分本文直接证据、前人研究及作者推测。图内英文若保留，需附完整中文对照。
\section*{文章信息}作者单位、贡献、利益冲突、基金、资助方角色、致谢及数据共享均按原文核对。样板不提供真实声明或基金编号。
\SampleBibliography
\SampleSourceInfo
\end{multicols}
\end{document}
